\section{A relative-zeta envelope for the remaining real zero}
\label{sec:s5-relative-variational}

The positive completed divisor of a relative Dedekind zeta
function removes the base field's entire zero contribution.
This gives one finite-family refinement of MA-RB-008.
It does not assert that the remaining interval is empty,
and does not assume noncoincidence of zeros of different
Artin factors.

\paragraph{Review status.}
The complete post-writing author reading and two independent
whole-proof readings have passed. The primary-source checks
have their own exact recorded scope.

Keep the actual Calegari extension
\(L/\mathbb Q\), \(G=\operatorname{Gal}(L/\mathbb Q)=S_5\),
and the irreducible representations
\(\mathbf1,\eta,V,V\eta,W,W\eta,X\) of MA-RB-002.
For every nontrivial subgroup \(H\le G\), set
\[
 \begin{gathered}
 K_H=L^H,\qquad n_H=[K_H:\mathbb Q]=[G:H],\qquad
 d_H=120-n_H,\\
 h_U(H)=\dim U^H,\qquad
 m_H=5-h_W(H)
        +\min\{4-h_{V\eta}(H),\,6-h_X(H)\}.
 \end{gathered}
 \tag{MA-RB.170}
\]
Conjugate subgroups give the same data, so one may take
one representative of each conjugacy class. No subgroup
enumeration is needed below. Write \((r_{1,H},r_{2,H})\)
for the signature of \(K_H\), and put
\[
 \begin{gathered}
 a_H=60-r_{1,H}-r_{2,H},\qquad b_H=60-r_{2,H},\\
 D_H(s)=\frac{\zeta_L(s)}{\zeta_{K_H}(s)},\qquad
 \Xi_H(s)=
 \left(\frac{\Delta_L}{\Delta_{K_H}}\right)^{s/2}
       \Gamma_{\mathbb R}(s)^{a_H}
       \Gamma_{\mathbb R}(s+1)^{b_H}D_H(s).
 \end{gathered}
 \tag{MA-RB.171}
\]
Here \(\Gamma_{\mathbb R}(s)=\pi^{-s/2}\Gamma(s/2)\),
and \(\psi_{\mathbb R}=\Gamma_{\mathbb R}'/
\Gamma_{\mathbb R}\). Since \(H\ne\{1\}\),
\(n_H\le60\); in particular \(a_H,b_H\ge0\)
and \(a_H+b_H=d_H>0\).

\paragraph{MA-RB-012 (a strictly smaller relative-zeta gate).}
Every \(D_H\) is entire. Its completion \(\Xi_H\)
is entire of order at most one, has nonzero values at
zero and one, and satisfies
\(\Xi_H(s)=\Xi_H(1-s)\) and real conjugation symmetry.
Its zero divisor is nonempty and is contained, with
nonnegative integer multiplicities, in the nontrivial
completed divisor of \(\zeta_L\).

For real \(\sigma>1\), define the exact expression
\[
 \begin{split}
 \mathscr H_H(\sigma)
   ={}&\frac12\log\frac{\Delta_L}{\Delta_{K_H}}
       +a_H\psi_{\mathbb R}(\sigma)
       +b_H\psi_{\mathbb R}(\sigma+1)\\
      &+\frac{\zeta_L'}{\zeta_L}(\sigma)
       -\frac{\zeta_{K_H}'}{\zeta_{K_H}}(\sigma).
 \end{split}
 \tag{MA-RB.172}
\]
There are no pole terms in this formula. It extends
continuously to \(\sigma=1\) and is strictly positive
on \([1,\infty)\). For \(k\ge1\), let
\[
 \begin{split}
 F_{H,k}(\sigma)
    &=(\sigma-\tfrac12)^2
        -\frac{m_Hk(2\sigma-1)}{\mathscr H_H(\sigma)},\\
 J_{H,k}&=\inf_{\sigma>1}F_{H,k}(\sigma),\qquad
 \mathfrak J_{L,k}
      =\min_{\{1\}\ne H\le G}J_{H,k}.
 \end{split}
 \tag{MA-RB.173}
\]
Each infimum is finite and is a minimum after adjoining
the endpoint \(\sigma=1\). It need not be attained at
an interior point. One has, unconditionally,
\[
          \mathfrak J_{L,k}
                 <\mathcal J_{\Delta_L,k}<\frac14,
 \tag{MA-RB.174}
\]
where \(\mathcal J_{\Delta_L,k}\) is the previously
checked discriminant-aware bound of MA-RB-008.

If the actual comparison-failure branch of MA-RB-002
occurs at its real abscissa \(\gamma\), with pole order
\(k\ge1\), then
\[
 \begin{gathered}
 \operatorname{ord}_{\gamma}\Xi_H
   =\operatorname{ord}_{1-\gamma}\Xi_H\ge m_Hk
                                      \quad(H\ne\{1\}),\\
             (\gamma-\tfrac12)^2\le\mathfrak J_{L,k}
                                      \le\mathfrak J_{L,1}.
 \end{gathered}
 \tag{MA-RB.175}
\]
Consequently \(\mathfrak J_{L,1}\le0\) excludes that
branch. Otherwise its sufficient shared-zero gate
is confined to
\[
 \begin{gathered}
 B(\gamma)=0,\qquad
        L(\gamma,V\eta)L(\gamma,X)=0,\\
 \frac12<\gamma
       \le\frac12+\sqrt{\mathfrak J_{L,1}}<1.
 \end{gathered}
 \tag{MA-RB.176}
\]
The new interval is strictly smaller than the nonempty
interval in MA-RB.90. These assertions use neither
\(\mathsf A\), \(M\), nor \((\mathcal R_B)\).

\begin{proof}
\emph{The classical divisibility input in its exact established form.}
Aramata's main theorem states that, for a finite normal
extension of number fields \(L/K\), the finite quotient
\(\zeta_L/\zeta_K\) is entire. The exact primary locator
is Aramata (1933), p.~31, the main statement preceding
Satz~1; its proof is completed on p.~34 following
Satz~6. All four pages of that proof were freshly
read in the source audit. This is a cited classical
input, not a new holomorphy theorem.
Here \(L/K_H\) is normal with Galois group \(H\),
even when \(K_H/\mathbb Q\) is not normal. Therefore
\[
             D_H=\zeta_L/\zeta_{K_H}\quad\hbox{is entire}.
 \tag{MA-RB.177}
\]
The statement concerns the full Dedekind functions,
including every ramified factor. It does not require
Artin holomorphy for individual irreducible factors.
The Artin induction and nontrivial Hecke holomorphy
interface is also explicitly recorded in the supplied
Hu--Kaneko--Martin--Schildkraut source,
\texttt{lem:artin-L}(a)--(d) and
\texttt{cor:hol-quotient}, including the latter's
complete proof. Tags MA-RB.178--.180 are reserved
and unused after this citation-only consolidation.

\emph{Completion and the genuine excess divisor.}
For a number field \(K\), use the same normalization
as MA-RB-001:
\[
 \xi_K(s)=s(s-1)\Delta_K^{s/2}
      \Gamma_{\mathbb R}(s)^{r_1(K)}
      \Gamma_{\mathbb C}(s)^{r_2(K)}\zeta_K(s),
 \qquad
 \Gamma_{\mathbb C}(s)
       =\Gamma_{\mathbb R}(s)\Gamma_{\mathbb R}(s+1).
 \tag{MA-RB.181}
\]
The standard Dedekind continuation, functional
equation and fixed-strip bounds used in MA-RB.8--.10
give an entire function of order at most one,
nonzero at zero and one, with real conjugation
symmetry and \(\xi_K(s)=\xi_K(1-s)\).
Its zeros lie in \(0\le\operatorname{Re}s\le1\);
their count in \(|s|\le T\) is \(O_K(T\log(T+3))\).
The proof of MA-RB.8--.10 uses these properties
for a general Dedekind completion; the particular
signature there enters only its explicit gamma
derivative.

Complex conjugation in \(G\) has type \(2B\), so
the signature of \(L\) is \((0,60)\).
Taking the quotient of (MA-RB.181) for \(L\) and
\(K_H\) cancels both factors \(s(s-1)\) and gives
exactly \(\Xi_H\) in (MA-RB.171).
At a nontrivial completed zero of \(\xi_{K_H}\)
the archimedean gamma factors are holomorphic and
nonzero: the only possible gamma endpoint in
the containing strip is zero, and \(\xi_{K_H}(0)\ne0\).
Entireness of \(D_H\) consequently says that
the zero occurs in \(\xi_L\) with at least the
same order. There are no other zeros of
\(\xi_{K_H}\). The quotient
\[
                   \Xi_H=\xi_L/\xi_{K_H}
 \tag{MA-RB.182}
\]
is therefore entire, with a nonnegative excess
divisor \(\mathcal Z_H\), and is nonzero at both
endpoints. This also verifies cancellation of
every possible trivial-zero or gamma singularity
in the formula (MA-RB.171); no hidden endpoint
pole or zero is retained.

For completeness its order can be checked without
an estimate for a reciprocal near zeros.
Write the genus-one Hadamard products for the
two functions \(\xi_L,\xi_{K_H}\), with
\(E_1(z)=(1-z)e^z\).
The factors belonging to \(\xi_{K_H}\) cancel
exactly, because the products converge locally
normally and \(\sum_\rho|\rho|^{-2}<\infty\).
Thus
\[
 \Xi_H(s)=\Xi_H(0)e^{\alpha_Hs}
                     \prod_{\rho\in\mathcal Z_H}E_1(s/\rho),
                       \qquad \alpha_H\in\mathbb R.
 \tag{MA-RB.183}
\]
The excess count is bounded by that of \(\xi_L\).
For \(|s|\le R\), the factors with \(|\rho|\le2R\)
have total log-modulus at most \(O_L(R\log^2(R+3))\):
use
\(\log|E_1(s/\rho)|
 \le\log(1+R/|\rho|)+R/|\rho|\)
and partial summation of the zero count.
For \(|\rho|>2R\), use
\(\log|E_1(s/\rho)|\le C R^2/|\rho|^2\);
their total is \(O_L(R\log(R+3))\).
The finite minimum nonzero zero modulus is harmless.
This proves order at most one.
Reflection and conjugation follow directly from
(MA-RB.182).

\emph{The exact positive logarithmic derivative.}
In (MA-RB.183),
\(\sum_{\rho\in\mathcal Z_H}\operatorname{Re}(1/\rho)\)
converges absolutely by the containing strip.
Call this sum \(S_H\). Logarithmic differentiation
at zero and one and the reflection
\(\rho\mapsto1-\rho\) give
\(-\alpha_H=\alpha_H+2S_H\), hence
\(\alpha_H=-S_H\).
For every real \(\sigma>1\) this fixes the constant
in the Hadamard formula and proves
\[
 \frac{\Xi_H'}{\Xi_H}(\sigma)
      =\sum_{\rho\in\mathcal Z_H}
                    \operatorname{Re}\frac1{\sigma-\rho}
      =\mathscr H_H(\sigma).
 \tag{MA-RB.184}
\]
The displayed real-part series is absolutely
convergent. The second equality is the direct
logarithmic derivative of (MA-RB.171), so in
particular no \(1/\sigma\) or \(1/(\sigma-1)\)
term is present.

The divisor \(\mathcal Z_H\) cannot be empty.
The gamma estimate already proved from the
gamma integral in MA-RB-007 gives
\(\psi_{\mathbb R}(\sigma+e)
 =\tfrac12\log\sigma-\tfrac12\log(2\pi)
     +O(1/\sigma)\) for \(e=0,1\).
Both finite Dedekind logarithmic derivatives
tend to zero on the positive real axis, by
dominated convergence in their absolutely
convergent Euler series at \(\sigma\ge2\).
Consequently
\[
          \mathscr H_H(\sigma)
               =\frac{d_H}{2}\log\sigma+O_L(1)
                                      \quad(\sigma\longrightarrow\infty).
 \tag{MA-RB.185}
\]
An empty divisor would make the first expression
in (MA-RB.184) identically zero, contradicting
\(d_H>0\). Every term in (MA-RB.184) is positive
when \(\sigma>1\), which proves strict positivity.

The function \(\Xi_H\) is nonzero at one, so its
logarithmic derivative extends continuously there.
The real-part zero series also extends to one:
its tail is locally bounded by a constant times
\(\sum_{|\rho|>2}|\rho|^{-2}\), and the remaining
finite terms have no singularity at one.
At one all terms are nonnegative. The reflection
symmetry ensures that a nonempty divisor contains
at least one zero with real part strictly less
than one, even if some zeros lie on a boundary
line. That zero contributes strictly positively
at one. Therefore \(\mathscr H_H(1)>0\).

\emph{The residual multiplicity without an extra zero-location premise.}
The permutation representation of the embeddings of
\(K_H\) is \(\operatorname{Ind}_H^G\mathbf1\).
Frobenius reciprocity gives multiplicity \(h_U(H)\)
for each irreducible \(U\). Its all-prime
factorization and the regular factorization imply
\[
      D_H(s)=
           \prod_U L(s,U)^{\,\dim U-h_U(H)}.
 \tag{MA-RB.186}
\]
All exponents are nonnegative integers.
At the particular residual point \(\gamma\),
MA-RB.19--.26 proves the following pointwise
orders without assuming general holomorphy of
the unknown five-dimensional factors:
\[
 \begin{gathered}
 x_{\mathbf1}=0,\quad x_\eta=b\ge0,\quad
 x_V=v_0\ge0,\quad x_{V\eta}=v_1\ge0,\quad
 x_X=e\ge0,\\
 x_W=n+k\ge k,\qquad x_{W\eta}\ge0,\qquad
                         v_1+e\ge k.
 \end{gathered}
 \tag{MA-RB.187}
\]
In particular \(x_{W\eta}\ge0\) follows there
from \(L_{W\eta}=\zeta_F/\zeta_{\mathbb Q}\)
at the real point \(\gamma\in(1/2,1)\), where
\(\zeta_{\mathbb Q}\ne0\).
Taking orders in (MA-RB.186), discarding only
nonnegative summands, and using \(v_1+e\ge k\)
gives
\[
 \begin{split}
 \operatorname{ord}_\gamma D_H
   &\ge (5-h_W)(n+k)
         +(4-h_{V\eta})v_1+(6-h_X)e\\
   &\ge m_H k .
 \end{split}
 \tag{MA-RB.188}
\]
The gamma factors are units at \(\gamma\), so
this is also the order of \(\Xi_H\).
Its functional equation gives the same order
at \(1-\gamma\). The two real points are distinct.
Moreover \(m_H\ge1\): the character row of \(W\)
has value strictly less than five at every
nonidentity element, so \(W\) is faithful.
For nontrivial \(H\), \(h_W(H)<5\); the other
two terms defining the minimum are nonnegative.

Retain these two real zeros in the positive sum
(MA-RB.184). With \(u=\gamma-1/2\) and
\(x=\sigma-1/2>1/2\), one obtains
\[
      \frac{2m_H kx}{x^2-u^2}
                      \le\mathscr H_H(\sigma),
       \qquad u^2\le F_{H,k}(\sigma).
 \tag{MA-RB.189}
\]
All denominators are positive because
\(0<u<1/2<x\).
Taking the infimum and then the finite minimum
proves (MA-RB.175), apart from its final
monotonicity assertion.

\emph{The minima and the strict comparison with the old gate.}
Strict positivity on \([1,\infty)\) makes every
\(F_{H,k}\) continuous there. Formula (MA-RB.185)
gives \(F_{H,k}(\sigma)\to+\infty\) as
\(\sigma\to\infty\), while
\[
             F_{H,k}(1)
                    =\frac14-\frac{m_Hk}{\mathscr H_H(1)}
 \tag{MA-RB.190}
\]
is finite. Hence the infimum over the open
half-line equals a minimum on \([1,\infty)\).
This proves the asserted finiteness and the
endpoint qualification.

Take \(H=A_5\), so \(K_H=E\) is the real quadratic
subfield. Its reality follows because complex
conjugation \(2B\) is even. The only irreducible
representations with \(A_5\)-invariants are
\(\mathbf1,\eta\): an invariant subspace for a
normal subgroup is \(G\)-stable, and an irreducible
representation trivial on \(A_5\) factors through
\(G/A_5=C_2\). Thus \(m_{A_5}=9\).

Write \(\mathscr H_K^{\rm full}=\xi_K'/\xi_K\)
on \(\sigma>1\). The same exact Hadamard
argument as (MA-RB.184) shows
\(\mathscr H_E^{\rm full}(\sigma)>0\).
Its divisor is nonempty: otherwise its zero
sum would vanish identically, whereas the
degree-two gamma derivative gives
\(\mathscr H_E^{\rm full}(\sigma)
 =\log\sigma+O_E(1)\) at infinity.
Now (MA-RB.182), followed by the exact positive
Euler-remainder identity MA-RB.95, gives
\[
 0<\mathscr H_{A_5}(\sigma)
    =\mathscr H_L^{\rm full}(\sigma)
                    -\mathscr H_E^{\rm full}(\sigma)
    <\mathscr H_L^{\rm full}(\sigma)
    \le\mathcal H_{\Delta_L}(\sigma).
 \tag{MA-RB.191}
\]
Let \(\sigma_*>1\) be an attained minimizer for
\(\mathcal J_{\Delta_L,k}\), whose existence
was proved in MA-RB-008. Since the numerator
\(9k(2\sigma_*-1)\) is strictly positive,
(MA-RB.191) implies
\[
 \mathfrak J_{L,k}\le J_{A_5,k}
    \le F_{A_5,k}(\sigma_*)
    <\mathcal F_{\Delta_L,k}(\sigma_*)
     =\mathcal J_{\Delta_L,k}.
 \tag{MA-RB.192}
\]
This proves (MA-RB.174). The strict step is
available at an old interior minimizer, so it
does not require an interior minimizer for
the new function.

There is no need for \(\mathsf A\) even in this
quadratic case. Indeed \(\zeta_E=\zeta_{\mathbb Q}C\)
has order \(b\) at \(\gamma\), and MA-RB.26
already gives
\[
       \operatorname{ord}_\gamma D_{A_5}
          =\operatorname{ord}_\gamma\zeta_L-b
                              \ge5n+9k.
 \tag{MA-RB.193}
\]
The possible quadratic zero is subtracted
exactly, including its multiplicity.
Finally each \(F_{H,k}\) decreases with \(k\),
as \(m_H>0\); the infima and finite minimum
do also. Together with MA-RB-002's two vanishing
factors, this proves (MA-RB.175)--(MA-RB.176)
and the nonpositive-endpoint exclusion.
\end{proof}

\paragraph{Three directly verified subgroup choices.}
The following examples are part of the same finite
envelope, not successive new reductions:
\[
\begin{array}{c|r|rr|rrr|r|rr}
 H&n_H&r_{1,H}&r_{2,H}&h_W&h_{V\eta}&h_X&m_H&a_H&b_H\\ \hline
 A_5&2&2&0&0&0&0&9&58&60\\
 \langle(12)\rangle&60&0&30&3&1&3&5&30&30\\
 \langle(12345)\rangle&24&0&12&1&0&2&8&48&48
\end{array}
 \tag{MA-RB.194}
\]
For the transposition, the invariant dimensions
are the averages of the identity and \(2A\)
character values. The relevant \(2A\) values
are \(1,-2,0\). For a five-cycle subgroup,
its four nonidentity elements all have type \(5A\);
the relevant values are \(0,-1,1\).
These give exactly the displayed averages.
In the last two cases the subgroup contains
no element of type \(2B\). Complex conjugation
therefore fixes no coset in \(G/H\), because
a fixed coset would put a conjugate of \(2B\)
inside \(H\). Thus these fields are totally
imaginary, giving the displayed signatures.
This uses the existing checked character rows;
it is not a new finite-representation certificate.

\paragraph{A finite Euler evaluation with an explicit error.}
For clarity this gate uses genuine field data,
not an assigned value for an unknown zero sum.
For a number field \(K\), define its rational
prime-power coefficients by
\[
 -\frac{\zeta_K'}{\zeta_K}(s)
       =\sum_{n\ge2}\frac{\Lambda_K(n)}{n^s},
 \qquad
 \Lambda_K(p^j)
       =(\log p)\sum_{\substack{\mathfrak p\mid p\\
                               f_{\mathfrak p}\mid j}}
                                     f_{\mathfrak p}.
 \tag{MA-RB.195}
\]
The factorization of every rational prime gives
these coefficients, including ramified primes.
They are zero at integers that are not prime powers.
The displayed sum is at most \([K:\mathbb Q]\log p\),
since \(\sum_{\mathfrak p\mid p}e_{\mathfrak p}
 f_{\mathfrak p}=[K:\mathbb Q]\).
Put \(c_H(n)=\Lambda_L(n)-\Lambda_{K_H}(n)\).
One always has the sufficient elementary bound
\[
              |c_H(n)|\le(120+n_H)\log n .
 \tag{MA-RB.196}
\]
In particular its constant for \(H=A_5\) is
\(122\). No sign of \(c_H(n)\) is assumed.
For an integer \(N\ge3\) and \(\sigma>1\), set
\[
 \begin{split}
 \mathscr H_{H,N}(\sigma)
    ={}&\frac12\log\frac{\Delta_L}{\Delta_{K_H}}
       +a_H\psi_{\mathbb R}(\sigma)
       +b_H\psi_{\mathbb R}(\sigma+1)
       -\sum_{2\le n\le N}c_H(n)n^{-\sigma},\\
 T_{H,N}(\sigma)
    ={}&(120+n_H)N^{1-\sigma}
       \left(\frac{\log N}{\sigma-1}
                         +\frac1{(\sigma-1)^2}\right).
 \end{split}
 \tag{MA-RB.197}
\]
The function \((\log x)x^{-\sigma}\) is decreasing
for \(x\ge3\). Comparison of the absolute tail
with its integral on \([N,\infty)\) proves
\[
 \left|\mathscr H_H(\sigma)-\mathscr H_{H,N}(\sigma)\right|
             \le T_{H,N}(\sigma),\qquad
 U_{H,N}(\sigma):=\mathscr H_{H,N}(\sigma)+T_{H,N}(\sigma)
                    \ge\mathscr H_H(\sigma)>0.
 \tag{MA-RB.198}
\]
Therefore a comparison failure would satisfy
the fully finite upper bound
\[
 (\gamma-\tfrac12)^2
      \le(\sigma-\tfrac12)^2
               -\frac{m_Hk(2\sigma-1)}{U_{H,N}(\sigma)}.
 \tag{MA-RB.199}
\]
For every fixed \(\sigma>1\), its right side
converges to \(F_{H,k}(\sigma)\) as \(N\to\infty\).
Taking rational \(\sigma>1\), finite \(N\), and
the finite subgroup family consequently recovers
\(\mathfrak J_{L,k}\) as the infimum of these
certified upper bounds. To see equality, the
upper bounds are all at least the true infimum;
conversely choose a real \(\sigma>1\) arbitrarily
close to an infimum, approximate it by a rational
one using continuity, and then take \(N\) large.
This argument also covers a minimum at one by
approaching from the right. It does not use an
Euler series at \(\sigma=1\), or assert a uniform
tail bound there.

\paragraph{Scope of positive products of these relative factors.}
The finite envelope already contains the bound from
any positive integral product of the same \(D_H\)
when only the guaranteed multiplicities \(m_Hk\)
are retained. Indeed, for nonnegative coefficients
\(u_H\), not all zero, and each \(\sigma>1\),
\[
 \frac{\sum_Hu_Hm_H}{\sum_Hu_H\mathscr H_H(\sigma)}
 =\sum_H
   \frac{u_H\mathscr H_H(\sigma)}
        {\sum_Ju_J\mathscr H_J(\sigma)}
             \frac{m_H}{\mathscr H_H(\sigma)}
 \le\max_H\frac{m_H}{\mathscr H_H(\sigma)}.
\]
For integer \(u_H\), the numerator is the guaranteed
multiplicity divided by \(k\), and the denominator is
the logarithmic derivative of the completed product.
Its corresponding two-real-zero upper bound is
therefore at least \(\min_H F_{H,k}(\sigma)\).
Taking an infimum in \(\sigma\) cannot lower it
below \(\mathfrak J_{L,k}\), since this family is
finite. The pointwise algebra also holds for real
nonnegative coefficients, without defining fractional
holomorphic powers. Higher actual multiplicities or
additional arithmetic information are outside this
limited optimality assertion.

\paragraph{Finite-cover and source scope.}
There is a genuine \(G\)-equivariant inclusion
\(\mathbb Q[G/H]\hookrightarrow\mathbb Q[G]\),
given by \(gH\mapsto\sum_{h\in H}gh\).
Right averaging
\(|H|^{-1}\sum_{h\in H}R_h\) is an idempotent
onto this image and commutes with the left \(G\)-action.
Its kernel is therefore a rational complement, an actual
finite Artin representation with character
\(\operatorname{Reg}_G-\operatorname{Ind}_H^G\mathbf1\);
equivalently it is the induction from \(H\) of
the augmentation representation of \(H\).
This describes the degree-zero-cohomology
Artin-motive quotient for the finite cover
\(\operatorname{Spec}L\to\operatorname{Spec}K_H\),
viewed over \(\mathbb Q\). Inertia invariants are
exact in characteristic zero, so the local
quotient includes every ramified Euler factor.
The conductor is \(\Delta_L/\Delta_{K_H}\) by
the same conductor--discriminant convention as
in the exact completion. Its Frobenius traces
need not be nonnegative. The positive measure
in the proof is the completed excess zero
divisor supplied by holomorphic divisibility,
not coefficientwise Euler positivity.

The primary divisibility source is Aramata
(1933), \emph{Über die Teilbarkeit der Dedekindschen
Zetafunktionen}, \emph{Proceedings of the Imperial
Academy}, \(9(2)\), 31--34,
doi:10.3792/pia/1195580866:
the principal statement is on p.~31, its
finite-character construction is in Sätze~1--6,
and the final analytic proof is on p.~34.
All four primary pages were freshly read in
their rendered form. The paper invokes this
known theorem in its exact normal-extension form,
as (MA-RB.177), without reproducing its proof.
The supplied Hu--Kaneko--Martin--Schildkraut
primary LaTeX, \texttt{HigherOrderZeros.tex},
\texttt{lem:artin-L}(a)--(d),
\texttt{lem:aramata-brauer}, and
\texttt{cor:hol-quotient}, was also read for
its exact all-prime induction/holomorphy scope.
No general Artin holomorphy theorem is substituted
for the relative divisibility statement.

The Dedekind continuation/growth inputs and
fixed Hadamard constant are the established
NT-Z10 and MA-RB.8--.12; the actual residual
orders are MA-RB.19--.26. The strict prior
denominator comparison uses MA-RB.95 and the
attained old minimum in MA-RB.87--.88.
For the categorical interpretation the inherited
exact locators remain SGA~1, Exposé~V, \S7,
and Deligne (1979), \emph{Valeurs de fonctions
\(L\) et périodes d'intégrales},
Definition~6.1, \S\S6.2--6.3 and equation~(0.5.3).
The conductor--discriminant locator remains
Neukirch, VII(11.7)--(11.9), printed pp.~533--534.
These previously checked categorical and
conductor locators are not claimed to have
received a fresh complete-source audit here.

This is a strictly smaller necessary interval
for the actual failure branch and hence a
weaker sufficient shared-zero exclusion
condition. It is not a proof that every
\(\mathfrak J_{L,1}\) is nonpositive, a proof
of \((\mathcal R_B)\), or a construction of
an arithmetic positive-excess example.
